% Options for packages loaded elsewhere
\PassOptionsToPackage{unicode}{hyperref}
\PassOptionsToPackage{hyphens}{url}
\PassOptionsToPackage{dvipsnames,svgnames,x11names}{xcolor}
%
\documentclass[
  11pt,
]{article}
\usepackage{amsmath,amssymb}
\usepackage{setspace}
\usepackage{iftex}
\ifPDFTeX
  \usepackage[T1]{fontenc}
  \usepackage[utf8]{inputenc}
  \usepackage{textcomp} % provide euro and other symbols
\else % if luatex or xetex
  \usepackage{unicode-math} % this also loads fontspec
  \defaultfontfeatures{Scale=MatchLowercase}
  \defaultfontfeatures[\rmfamily]{Ligatures=TeX,Scale=1}
\fi
\usepackage{lmodern}
\ifPDFTeX\else
  % xetex/luatex font selection
  \setmainfont[]{DejaVu Serif}
  \setmathfont[]{DejaVu Math TeX Gyre}
\fi
% Use upquote if available, for straight quotes in verbatim environments
\IfFileExists{upquote.sty}{\usepackage{upquote}}{}
\IfFileExists{microtype.sty}{% use microtype if available
  \usepackage[]{microtype}
  \UseMicrotypeSet[protrusion]{basicmath} % disable protrusion for tt fonts
}{}
\makeatletter
\@ifundefined{KOMAClassName}{% if non-KOMA class
  \IfFileExists{parskip.sty}{%
    \usepackage{parskip}
  }{% else
    \setlength{\parindent}{0pt}
    \setlength{\parskip}{6pt plus 2pt minus 1pt}}
}{% if KOMA class
  \KOMAoptions{parskip=half}}
\makeatother
\usepackage{xcolor}
\usepackage[margin=1in]{geometry}
\setlength{\emergencystretch}{3em} % prevent overfull lines
\providecommand{\tightlist}{%
  \setlength{\itemsep}{0pt}\setlength{\parskip}{0pt}}
\setcounter{secnumdepth}{-\maxdimen} % remove section numbering
\usepackage{amsmath,amssymb,mathtools}
\usepackage{enumitem}
\setlist[itemize]{topsep=3pt,itemsep=2pt,parsep=0pt}
\setlist[enumerate]{topsep=3pt,itemsep=2pt,parsep=0pt}
\usepackage{microtype}
\usepackage{titlesec}
\titleformat{\section}{\large\bfseries}{\thesection}{0.7em}{}
\titleformat{\subsection}{\normalsize\bfseries}{\thesubsection}{0.7em}{}
\usepackage[most]{tcolorbox}
\tcbset{colback=gray!4,colframe=gray!40,arc=2pt,boxrule=0.4pt,left=6pt,right=6pt,top=5pt,bottom=5pt}
\ifLuaTeX
  \usepackage{selnolig}  % disable illegal ligatures
\fi
\usepackage{bookmark}
\IfFileExists{xurl.sty}{\usepackage{xurl}}{} % add URL line breaks if available
\urlstyle{same}
\hypersetup{
  pdftitle={Recommendation on CPSV16 Step 1, One-Copy BNT, and Non-Dominant Projection},
  colorlinks=true,
  linkcolor={blue},
  filecolor={Maroon},
  citecolor={Blue},
  urlcolor={blue},
  pdfcreator={LaTeX via pandoc}}

\title{Recommendation on CPSV16 Step 1, One-Copy BNT, and Non-Dominant
Projection}
\author{}
\date{May 13, 2026}

\begin{document}
\maketitle

\setstretch{1.08}
\section{Recommendation on CPSV16 Step 1, one-copy BNT, and non-dominant
projection}\label{recommendation-on-cpsv16-step-1-one-copy-bnt-and-non-dominant-projection}

\subsection{Executive summary}\label{executive-summary}

The current issue is not merely a missing estimate. It is a mismatch
between three things:

\begin{enumerate}
\def\labelenumi{\arabic{enumi}.}
\tightlist
\item
  the paper-level BNT/canonical-form structure;
\item
  the current formalization surface, apparently called
  \texttt{IsCanonicalFormBNT};
\item
  the proof strategy being attempted for CPSV16 Step 1.
\end{enumerate}

The current \texttt{IsCanonicalFormBNT} surface is described as carrying
a single weight per sector with strict modulus ordering. This is a
\textbf{one-copy specialization} of the paper's two-layer structure,
where each BNT sector may contain several gauge-phase-equivalent copies
and contributes a power-sum coefficient.

The non-dominant per-block projection argument should not be pushed
further. It works for the dominant block, but for a non-dominant block
both sides of the projected identity can decay to zero. The paper's Step
1 mechanism is instead linear independence of the relevant combined
family and coefficient comparison.

The examples below are meant to make the distinction concrete.

\begin{center}\rule{0.5\linewidth}{0.5pt}\end{center}

\subsection{1. Paper-level BNT structure is not
one-copy}\label{paper-level-bnt-structure-is-not-one-copy}

In the MPDO/CPGSV canonical-form framework, after grouping
gauge-phase-equivalent normal blocks into a basis of normal tensors, one
sector may contain several copies:

\[
A^i \simeq
\bigoplus_{j=1}^g
\bigoplus_{q=1}^{r_j}
\mu_{j,q} X_{j,q} A_j^i X_{j,q}^{-1}.
\]

At the MPV level this gives

\[
V_N(A)
=
\sum_{j=1}^g
\left(\sum_{q=1}^{r_j}\mu_{j,q}^N\right)V_N(A_j).
\]

So the coefficient attached to a BNT sector is generally a power sum

\[
c_j(N)=\sum_q \mu_{j,q}^N,
\]

not a single scalar power \(\mu_j^N\).

The uploaded note explicitly describes the current
\texttt{IsCanonicalFormBNT} surface as the specialization \(r_j=1\): one
weight per sector, strictly ordered by modulus. That is a genuine
restriction relative to the paper's BNT structure.

\begin{center}\rule{0.5\linewidth}{0.5pt}\end{center}

\subsection{2. Example: one BNT sector with two phase
copies}\label{example-one-bnt-sector-with-two-phase-copies}

Take one normal tensor \(C\), and consider the block tensor

\[
A = C \oplus (-C).
\]

Equivalently, this is one paper-level BNT tensor \(C\) with two scalar
copies

\[
\mu_1=1,\qquad \mu_2=-1.
\]

Then

\[
V_N(A)=1^N V_N(C)+(-1)^N V_N(C)
      =(1+(-1)^N)V_N(C).
\]

Thus

\[
V_N(A)=
\begin{cases}
2V_N(C), & N\text{ even},\\
0, & N\text{ odd}.
\end{cases}
\]

This is exactly the kind of phenomenon the paper's BNT-sector notation
records: repeated gauge-phase-equivalent copies of the same normal block
combine into one sector coefficient.

But it cannot be represented by a \textbf{single} scalar power
\(\mu^N V_N(C)\), because no scalar \(\mu\) satisfies

\[
\mu^N = 1+(-1)^N
\]

for all \(N\). The problem is not notation; the one-copy surface has
removed genuine multiplicity data.

This example should not be described as two distinct BNT elements. It is
one BNT element with multiplicity two. That is precisely the distinction
the current one-copy predicate appears to collapse.

\begin{center}\rule{0.5\linewidth}{0.5pt}\end{center}

\subsection{3. Example: equal-modulus
copies}\label{example-equal-modulus-copies}

Similarly, take

\[
A = C \oplus e^{i\theta}C.
\]

Then

\[
V_N(A)=(1+e^{iN\theta})V_N(C).
\]

Both scalar copies have modulus \(1\). A strict block-level ordering

\[
|\mu_1|>|\mu_2|>\cdots
\]

cannot keep both copies as separate ordered blocks. The paper handles
this by grouping them into one BNT sector and recording the coefficient

\[
1+e^{iN\theta}.
\]

Therefore, a strict one-weight-per-sector formalization is weaker than
the paper unless there is an additional theorem that reconstructs the
full sector/multiplicity structure.

\begin{center}\rule{0.5\linewidth}{0.5pt}\end{center}

\subsection{4. Why the non-dominant projection argument
fails}\label{why-the-non-dominant-projection-argument-fails}

The uploaded note's analytic point is correct. Suppose, after dominant
normalization, the assembled \(B\)-side has two blocks

\[
B_{\mathrm{tot}}=B_0\oplus \lambda B_1,
\qquad 0<|\lambda|<1.
\]

Assume the proportionality scalar has the dominant behavior

\[
c_N\sim 1.
\]

Project onto the dominant block \(B_0\). The leading RHS contribution is

\[
c_N\cdot 1^N\cdot
\langle V_N(B_0),V_N(B_0)\rangle
\to \ell_0\neq 0.
\]

So if the LHS tends to zero, this gives a contradiction.

But now project onto the non-dominant block \(B_1\). The leading RHS
contribution is

\[
c_N\lambda^N
\langle V_N(B_1),V_N(B_1)\rangle
\sim \lambda^N\ell_1
\to 0.
\]

Thus both sides of the projected identity may tend to zero. There is no
contradiction.

This is not a Lean technicality. It is a genuine analytic obstruction to
extending the dominant-block projection proof to arbitrary blocks.

\begin{center}\rule{0.5\linewidth}{0.5pt}\end{center}

\subsection{5. What the paper's Step 1 proof uses
instead}\label{what-the-papers-step-1-proof-uses-instead}

The source proof should be formalized as a
linear-independence/coefficient-comparison argument.

Start from eventual proportionality:

\[
\sum_j (\mu_j^A)^N V_N(A_j)
=
c_N\sum_k(\mu_k^B)^N V_N(B_k).
\]

Fix \(k_0\). Move all terms except \(B_{k_0}\) to the left:

\[
\sum_j (\mu_j^A)^N V_N(A_j)
-
c_N\sum_{k\neq k_0}(\mu_k^B)^N V_N(B_k)
=
c_N(\mu_{k_0}^B)^N V_N(B_{k_0}).
\]

If the actual family appearing in this relation,

\[
\{V_N(A_j)\}_j\cup \{V_N(B_k)\}_k,
\]

is eventually linearly independent, then coefficient comparison forces

\[
c_N(\mu_{k_0}^B)^N=0.
\]

This contradicts \(c_N\neq0\) and \(\mu_{k_0}^B\neq0\).

The crucial point is that this proof does \textbf{not} require

\[
c_N(\mu_{k_0}^B)^N
\]

to stay bounded away from zero. It only requires it to be nonzero. This
is why coefficient comparison works where projection estimates fail.

\begin{center}\rule{0.5\linewidth}{0.5pt}\end{center}

\subsection{6. Why the smaller linear-independence lemma is
insufficient}\label{why-the-smaller-linear-independence-lemma-is-insufficient}

It is not enough to know linear independence of

\[
\{V_N(A_j)\}_j\cup\{V_N(B_{k_0})\}.
\]

The exact vector relation also contains the other \(B_k\)'s:

\[
\sum_j a_j(N)V_N(A_j)
-
\sum_{k\neq k_0} b_k(N)V_N(B_k)
-
b_{k_0}(N)V_N(B_{k_0})=0.
\]

Independence of the smaller family gives no coefficient information in a
relation involving omitted vectors.

A finite-dimensional toy example shows the issue. In \(\mathbb C^2\),
let

\[
u=e_1,\qquad v=e_2,\qquad w=e_1+e_2.
\]

The pair \(\{u,v\}\) is linearly independent. But in the larger family,

\[
w-u-v=0.
\]

So knowing that \(\{u,v\}\) is independent does not let one isolate the
coefficient of \(v\) in a relation involving \(w\). Analogously, knowing
independence of \(\{A_j\}\cup\{B_{k_0}\}\) does not isolate \(B_{k_0}\)
in the proportionality identity involving all \(B_k\).

The needed statement is the full combined-family linear independence
lemma, or at least a version strong enough to isolate the coefficient of
\(B_{k_0}\) in the actual relation.

\begin{center}\rule{0.5\linewidth}{0.5pt}\end{center}

\subsection{7. Recommended message to the formalization
agents}\label{recommended-message-to-the-formalization-agents}

Please separate the following theorem surfaces explicitly:

\begin{enumerate}
\def\labelenumi{\arabic{enumi}.}
\item
  \textbf{One-copy BNT canonical form.}\\
  One scalar weight per BNT sector, strict ordering of moduli. This is
  the current \texttt{IsCanonicalFormBNT} surface as described in the
  uploaded note.
\item
  \textbf{Paper-level sector/BNT canonical form.}\\
  A sector may contain several gauge-phase-equivalent copies, and the
  coefficient of a BNT block is a power sum \(\sum_q\mu_{j,q}^N\).
\item
  \textbf{Fully unconditional canonical-form existence.}\\
  Starting from arbitrary tensors, one must still construct the
  appropriate canonical/BNT data before applying the uniqueness theorem.
\end{enumerate}

The present non-dominant Step 1 issue concerns the gap between (1) and
(2), plus an invalid attempt to prove a universal statement using a
dominant-block projection argument.

Please do not try to extend the dominant-block projection proof to
non-dominant blocks. For a non-dominant block, the scalar prefactor
decays, so the projected proportionality identity can reduce
asymptotically to \(0=0\). The two-block example
\(B_0\oplus\lambda B_1\), \(|\lambda|<1\), already shows this.

The correct CPSV/MPDO proof mechanism is eventual linear independence of
the full combined family and exact coefficient comparison.

\begin{center}\rule{0.5\linewidth}{0.5pt}\end{center}

\subsection{8. Recommended formalization
repairs}\label{recommended-formalization-repairs}

\begin{enumerate}
\def\labelenumi{\arabic{enumi}.}
\item
  Rename or document the current predicate as \textbf{one-copy BNT
  canonical form}.
\item
  Add or promote a paper-faithful sector structure:

  \[
  (\lambda_j,\{\omega_{j,q}\}_q,A_j),
  \qquad |\omega_{j,q}|=1,
  \]

  with sector coefficient

  \[
  \lambda_j^N\sum_q\omega_{j,q}^N.
  \]
\item
  Prove the full combined-family Lem1:

  \[
  \text{all relevant cross-overlaps decay}
  \Longrightarrow
  \{V_N(A_j)\}_j\cup\{V_N(B_k)\}_k
  \text{ is eventually linearly independent}.
  \]
\item
  Use exact coefficient comparison, not asymptotic projection, for
  arbitrary \(k_0\).
\item
  Keep the dominant projection theorem as a useful special lemma, but do
  not present it as the proof of the universal nondecay statement.
\item
  Do not claim the full MPDO/CPGSV BNT theorem from the one-copy
  predicate unless a sector/grouping theorem has been added.
\end{enumerate}

\begin{center}\rule{0.5\linewidth}{0.5pt}\end{center}

\subsection{Final classification}\label{final-classification}

The current situation is best classified as:

\begin{itemize}
\tightlist
\item
  \textbf{scope restriction:} the current predicate is one-copy and
  therefore weaker than the paper-level BNT structure;
\item
  \textbf{proof-strategy gap:} the non-dominant projection argument is
  analytically insufficient;
\item
  \textbf{repairable formalization gap:} the right repair is full
  combined-family linear independence plus coefficient comparison, or a
  move to the paper-faithful sector decomposition.
\end{itemize}

The examples above are deliberately simple. Their role is to prevent the
issue from being mistaken for a cosmetic naming complaint. The
formalization is not merely missing polish: it is using a restricted
canonical-form surface and a proof method that fails exactly where the
paper switches to BNT-sector linear independence.

\end{document}
